\clearpage
\phantomsection
\chapter{CONCLUSION}
\label{chap:conclusion}

\section{Summary of Achievements}

This thesis presented the design, implementation, and evaluation of an automated game-balancing system for a Roguelike Deckbuilder prototype, combining a Multi-Objective Evolutionary Algorithm (NSGA-II), Reinforcement Learning (RL), and a Deterministic Heuristic Agent (DHA). The system addresses one of the core challenges in modern game development: producing parameter configurations that are simultaneously balanced, engaging, and internally coherent, without requiring extensive manual tuning.

The following objectives were stated at the outset:

\begin{enumerate}
    \item \textbf{Build a complete, simulatable Roguelike Deckbuilder prototype.} \\ Achieved: The C\#/.NET 9 engine supports 30+ cards, 20+ relics, a 15-floor map with 5 room types, turn-based combat with 10+ status effects, and executes $\sim$2,500 complete game runs per second on a single core.

    \item \textbf{Design and implement a Deterministic Heuristic Agent for automated simulation.} \\ Achieved: The three-priority DHA (survival $>$ lethal $>$ optimal) runs complete game episodes without human intervention, serves as the primary evaluation oracle during optimization, and provides a stable baseline against which RL agents can be compared.

    \item \textbf{Train RL agents with distinct behavioral profiles.} \\ Achieved: Four MaskablePPO agents (Aggressive, Defensive, Balanced, Adaptive) were trained for 8 million steps each on Kaggle T4 GPUs. Win rates ranged from 24.5\% (Aggressive) to 34.0\% (Adaptive), covering a broad spectrum of play styles.

    \item \textbf{Optimize game parameters using a single-objective Genetic Algorithm.} \\ Achieved: The GA improved overall fitness from 1.45 to 32.66 ($+2,150\%$), win rate from 36.0\% to 38.3\%, and moved key metrics toward target values in 73 minutes.

    \item \textbf{Optimize using NSGA-II for multi-objective balance.} \\ Achieved: NSGA-II achieved Balance $= 0.960$, Engagement $= 0.900$, and Coherence $= 1.000$ (perfect, no trap cards), increased viable cards from 12 to 14, and ran 33\% faster than the GA (49 vs.\ 73 minutes).
\end{enumerate}

The headline results confirm that NSGA-II with a hierarchical genome is the superior approach. The 50-dimensional hierarchical genome is semantically meaningful (each dimension corresponds to an interpretable design concept such as ``early-game difficulty scaling''), converges 33\% faster, and achieves better outcomes on all measured dimensions simultaneously.

\section{Scientific and Practical Contributions}

\subsection{Scientific Contributions}

\begin{itemize}
    \item \textbf{Application of NSGA-II to Roguelike Deckbuilder balancing}: This work provides an end-to-end instantiation of the NSGA-II algorithm for the specific problem structure of card game balance, including the formalization of Balance, Engagement, and Coherence as objectives and the derivation of appropriate evaluation metrics from simulation data.

    \item \textbf{Hierarchical genome encoding for game parameter optimization}: The proposed HierarchicalGenome representation -- grouping the $\sim$500 primitive game parameters into $\sim$50 semantically coherent dimensions -- is a design principle applicable to any data-driven game that uses a JSON-based parameter system. It enables faster convergence and more interpretable mutations compared to flat parameter vectors.

    \item \textbf{Multi-behavioral RL agent ensemble for game evaluation}: Training four agents with distinct reward functions (Aggressive, Defensive, Balanced, Adaptive) and using them as an ensemble evaluator provides a richer signal of game quality than any single agent. The Adaptive agent's emergence of context-dependent strategy without explicit rule encoding demonstrates the strength of reward shaping over hand-coded heuristics for behavioral specification.

    \item \textbf{pythonnet as a high-performance simulation bridge}: The comparison of gRPC ($\sim$25\,ms/call) versus pythonnet ($<1\,\mu$s/call) provides an empirical argument for direct CLR embedding as the integration strategy of choice for .NET/Python hybrid AI pipelines.
\end{itemize}

\subsection{Practical Contributions}

\begin{itemize}
    \item A fully functional, open-source Roguelike Deckbuilder simulation engine in C\#/.NET 9, ready to be extended with additional content or adapted for research.
    \item A Python training pipeline runnable on free Kaggle notebooks -- no specialized hardware required.
    \item A modular optimization module (\texttt{GeneticAlgorithm.cs}, \texttt{NSGA2Optimizer.cs}, \texttt{HierarchicalGenome.cs}) that is game-engine-agnostic and can be applied to any game with a parameterizable JSON data layer.
    \item Empirical evidence that targeting a 30--35\% win rate (achievable by a well-trained RL agent) produces a game experience that is neither too easy nor too hard, which is a practically useful design guideline.
\end{itemize}

\section{Limitations}

Despite the positive results, the system has several important limitations that should be acknowledged.

\subsection{Agent Representativeness}

The DHA and all four RL agents are \textbf{computational proxies} for human players, not actual human players. The DHA plays optimally within its heuristic framework, which may differ systematically from human decision-making in edge cases (e.g., the DHA never bluffs, never makes emotional or exploratory sub-optimal plays, and never experiences decision fatigue). The RL agents learn from simulated experience and are trained against the game as-is; they may develop patterns that are effective but unusual or unintuitive to a human observer.

Consequently, a game configuration rated as ``balanced'' by this system should be considered a \textbf{necessary but not sufficient} condition for human-perceived balance. A human playtest study would be required to validate the outputs.

\subsection{Fixed Hierarchy Structure}

The hierarchical genome groups parameters into 50 dimensions based on the designer's semantic understanding of the game. If this grouping is incorrect (e.g., two parameters that seem related actually interact antagonistically), the NSGA-II will learn misleading patterns. The hierarchy was defined manually and has not been validated against alternative grouping schemes.

\subsection{Content Scope}

The prototype contains 30 cards, 20+ relics, and 15 floors -- sufficient for research validation but significantly smaller than commercial Roguelike Deckbuilders (e.g., \textit{Slay the Spire} has 75+ cards per character). It is unclear whether the optimization approach scales linearly or super-linearly with content size.

\subsection{Single-Character Evaluation}

All experiments use a single hero class. Multi-character games require cross-character balance (ensuring each character has roughly equal win rates) in addition to intra-character balance, introducing an additional objective that was not addressed here.

\section{Future Work}

The results of this thesis point to several promising directions for continuation.

\subsection{RL-Integrated Optimization}

The current system uses the DHA as the evaluation oracle during optimization (fast) and the RL agents only for validation (slow). A natural next step is to use the RL agents \textit{directly} in the optimization loop, possibly via a fast-evaluation approximation (e.g., a learned surrogate model trained on DHA+RL evaluation pairs). This would close the loop between the optimization and play-testing components, reducing the reliance on the DHA's heuristic approximation.

\subsection{NSGA-III and Many-Objective Extensions}

The current system optimizes three objectives. A richer balancing framework might include additional objectives such as:
\begin{itemize}
    \item Comeback potential (fraction of runs where a player at $<20\%$ HP recovers to win).
    \item Early-game vs. late-game content pacing.
    \item Cross-character win rate parity (for multi-character games).
\end{itemize}
With four or more objectives, the crowding distance mechanism of NSGA-II becomes less effective, and NSGA-III \cite{deb2014evolutionary} or decomposition-based methods (MOEA/D) would be more appropriate.

\subsection{Meta-Learning for Hierarchy Discovery}

Rather than manually designing the hierarchical genome, future work could use meta-learning or clustering analysis on the parameter correlation matrix (estimated from many random evaluations) to automatically discover the optimal parameter grouping. This would make the system fully autonomous and remove the dependency on designer domain knowledge.

\subsection{Human Playtest Integration}

Recording human playtest sessions and using them to train a player model (via inverse RL or behavioral cloning) would replace the proxy agents with a more faithful approximation of human play. This data could also be used to learn a reward function for RL training that directly targets human-perceived fun, rather than proxy metrics.

\subsection{Expanded Game Content}

Adding multiple hero classes, additional floor types (puzzles, boss variants), and a larger card pool would stress-test the scalability of the optimization approach and produce a prototype closer to commercial game quality.

\subsection{Designer Tool}

Packaging the optimization pipeline as an interactive designer tool -- with a web dashboard showing the Pareto front, card viability plots, and parameter sliders -- would make the system accessible to game designers without optimization expertise. The designer could drag points along the Pareto front and immediately see which specific parameter values correspond to their selected balance trade-off.
